\documentclass[10pt,a4paper]{amsart}
\usepackage{amsmath,amssymb,amsthm,graphicx}
% Retained minimal preamble for the Stacks Project chapter flat.tex (More on Flatness).
% Verbatim excerpt of the chapter preamble: Xy-pic setup (native preamble lines 14-18) and the native
% theorem declarations with the mathematical macros (native preamble lines 154-217). Copied mechanically;
% nothing added, removed or redefined. The wrapper class and the amsmath/amssymb/amsthm/graphicx base
% packages are supplied by the assembler, so the class/externaldocument part of the native preamble is omitted.
% For commutative diagrams we use Xy-pic
\usepackage[all]{xy}

% We use 2cell for 2-commutative diagrams.
\xyoption{2cell}
% Theorem environments.
%
\theoremstyle{plain}
\newtheorem{theorem}[subsection]{Theorem}
\newtheorem{proposition}[subsection]{Proposition}
\newtheorem{lemma}[subsection]{Lemma}

\theoremstyle{definition}
\newtheorem{definition}[subsection]{Definition}
\newtheorem{example}[subsection]{Example}
\newtheorem{exercise}[subsection]{Exercise}
\newtheorem{situation}[subsection]{Situation}

\theoremstyle{remark}
\newtheorem{remark}[subsection]{Remark}
\newtheorem{remarks}[subsection]{Remarks}

\numberwithin{equation}{subsection}

% Macros
%
\def\lim{\mathop{\mathrm{lim}}\nolimits}
\def\colim{\mathop{\mathrm{colim}}\nolimits}
\def\Spec{\mathop{\mathrm{Spec}}}
\def\Hom{\mathop{\mathrm{Hom}}\nolimits}
\def\Ext{\mathop{\mathrm{Ext}}\nolimits}
\def\SheafHom{\mathop{\mathcal{H}\!\mathit{om}}\nolimits}
\def\SheafExt{\mathop{\mathcal{E}\!\mathit{xt}}\nolimits}
\def\Sch{\mathit{Sch}}
\def\Mor{\mathop{\mathrm{Mor}}\nolimits}
\def\Ob{\mathop{\mathrm{Ob}}\nolimits}
\def\Sh{\mathop{\mathit{Sh}}\nolimits}
\def\NL{\mathop{N\!L}\nolimits}
\def\CH{\mathop{\mathrm{CH}}\nolimits}
\def\proetale{{pro\text{-}\acute{e}tale}}
\def\etale{{\acute{e}tale}}
\def\QCoh{\mathit{QCoh}}
\def\Ker{\mathop{\mathrm{Ker}}}
\def\Im{\mathop{\mathrm{Im}}}
\def\Coker{\mathop{\mathrm{Coker}}}
\def\Coim{\mathop{\mathrm{Coim}}}

% Boxtimes
%
\DeclareMathSymbol{\boxtimes}{\mathbin}{AMSa}{"02}

%
% Macros for moduli stacks/spaces
%
\def\QCohstack{\mathcal{QC}\!\mathit{oh}}
\def\Cohstack{\mathcal{C}\!\mathit{oh}}
\def\Spacesstack{\mathcal{S}\!\mathit{paces}}
\def\Quotfunctor{\mathrm{Quot}}
\def\Hilbfunctor{\mathrm{Hilb}}
\def\Curvesstack{\mathcal{C}\!\mathit{urves}}
\def\Polarizedstack{\mathcal{P}\!\mathit{olarized}}
\def\Complexesstack{\mathcal{C}\!\mathit{omplexes}}
% \Pic is the operator that assigns to X its picard group, usage \Pic(X)
% \Picardstack_{X/B} denotes the Picard stack of X over B
% \Picardfunctor_{X/B} denotes the Picard functor of X over B
\def\Pic{\mathop{\mathrm{Pic}}\nolimits}
\def\Picardstack{\mathcal{P}\!\mathit{ic}}
\def\Picardfunctor{\mathrm{Pic}}
\def\Deformationcategory{\mathcal{D}\!\mathit{ef}}

% --- appended 2026-10-12 (Chief): the project own definitions of the reference/slogan/history
% display environments, copied mechanically from the shared Stacks preamble lines 8-12
% (\usepackage{verbatim} plus the three \newenvironment lines, which the chapter preamble keeps
% verbatim from the project preamble). Nothing added, removed or redefined: these environments
% carry no mathematical content and the project itself typesets them as comments. Without these
% three lines a counted statement whose own text uses them fails with Environment slogan undefined.
\usepackage{verbatim}
\newenvironment{reference}{\comment}{\endcomment}
\newenvironment{slogan}{\comment}{\endcomment}
\newenvironment{history}{\comment}{\endcomment}
\usepackage{cleveref}
\crefname{originallemma}{original lemma}{original lemmas}
\Crefname{originallemma}{Original lemma}{Original lemmas}
\makeatletter
\newlabel{descent-lemma-descending-property-finite}{{02LA}{source}{}{sourcepointer-1}{}}
\newlabel{descent-lemma-descending-property-finite@cref}{{[originallemma][0][]02LA}{[1][0][]source}{}{}{}}
\makeatother
\crefname{originallemma}{original lemma}{original lemmas}
\Crefname{originallemma}{Original lemma}{Original lemmas}
\makeatletter
\newlabel{lemma-flat-finite-presentation-dimension-over-dense-open}{{081L}{source}{}{sourcepointer-2}{}}
\newlabel{lemma-flat-finite-presentation-dimension-over-dense-open@cref}{{[originallemma][0][]081L}{[1][0][]source}{}{}{}}
\makeatother
\crefname{originallemma}{original lemma}{original lemmas}
\Crefname{originallemma}{Original lemma}{Original lemmas}
\makeatletter
\newlabel{lemma-flat-finite-type-finitely-presented-over-dense-open-X}{{081P}{source}{}{sourcepointer-3}{}}
\newlabel{lemma-flat-finite-type-finitely-presented-over-dense-open-X@cref}{{[originallemma][0][]081P}{[1][0][]source}{}{}{}}
\makeatother
\crefname{originallemma}{original lemma}{original lemmas}
\Crefname{originallemma}{Original lemma}{Original lemmas}
\makeatletter
\newlabel{more-morphisms-lemma-etale-splits-off-quasi-finite-part}{{02LP}{source}{}{sourcepointer-4}{}}
\newlabel{more-morphisms-lemma-etale-splits-off-quasi-finite-part@cref}{{[originallemma][0][]02LP}{[1][0][]source}{}{}{}}
\makeatother
\crefname{originallemma}{original lemma}{original lemmas}
\Crefname{originallemma}{Original lemma}{Original lemmas}
\makeatletter
\newlabel{morphisms-lemma-finite-flat}{{02KB}{source}{}{sourcepointer-5}{}}
\newlabel{morphisms-lemma-finite-flat@cref}{{[originallemma][0][]02KB}{[1][0][]source}{}{}{}}
\makeatother
\crefname{originallemma}{original lemma}{original lemmas}
\Crefname{originallemma}{Original lemma}{Original lemmas}
\makeatletter
\newlabel{morphisms-lemma-finite-locally-free}{{04MH}{source}{}{sourcepointer-6}{}}
\newlabel{morphisms-lemma-finite-locally-free@cref}{{[originallemma][0][]04MH}{[1][0][]source}{}{}{}}
\makeatother
\crefname{originallemma}{original lemma}{original lemmas}
\Crefname{originallemma}{Original lemma}{Original lemmas}
\makeatletter
\newlabel{morphisms-lemma-flat-morphism-scheme-theoretically-dense-open}{{081H}{source}{}{sourcepointer-7}{}}
\newlabel{morphisms-lemma-flat-morphism-scheme-theoretically-dense-open@cref}{{[originallemma][0][]081H}{[1][0][]source}{}{}{}}
\makeatother
\crefname{originallemma}{original lemma}{original lemmas}
\Crefname{originallemma}{Original lemma}{Original lemmas}
\makeatletter
\newlabel{morphisms-lemma-fppf-open}{{01UA}{source}{}{sourcepointer-8}{}}
\newlabel{morphisms-lemma-fppf-open@cref}{{[originallemma][0][]01UA}{[1][0][]source}{}{}{}}
\makeatother
\crefname{originallemma}{original lemma}{original lemmas}
\Crefname{originallemma}{Original lemma}{Original lemmas}
\makeatletter
\newlabel{morphisms-lemma-locally-quasi-finite-rel-dimension-0}{{0397}{source}{}{sourcepointer-9}{}}
\newlabel{morphisms-lemma-locally-quasi-finite-rel-dimension-0@cref}{{[originallemma][0][]0397}{[1][0][]source}{}{}{}}
\makeatother
\title{A separated flat locally finite type morphism which is an isomorphism over a retrocompact scheme theoretically dense open is an open immersion}
\author{The Stacks Project Authors}
\date{stacks-project:flat.tex:local-snapshot-20261011; GNU FDL 1.2 or later}
\begin{document}
\maketitle
\noindent\textit{Editorial note.} All delivered mathematical text is the author's own, copied from the pinned Stacks Project chapter file flat.tex; no statement or proof step was rewritten, reordered or removed.\par\medskip
\noindent\textit{Editorial note.} Preamble provenance: the chapter's own retained minimal preamble (sources/stacks-flat/retained-preamble.tex, lines 1-85) is a verbatim excerpt of the native Xy-pic setup and of the native theorem declarations and mathematical macros, followed by the project's own reference/slogan/history display-environment definitions copied mechanically from the shared preamble; the wrapper class and the amsmath/amssymb/amsthm/graphicx base packages are supplied by the assembler, and extra\_packages is left empty because everything needed is already in the retained preamble. No macro was invented or redefined.\par\medskip
\noindent\textit{Editorial note.} Same-chapter pointers are located at their \textbackslash{}label lines inside the pinned chapter file; cross-chapter pointers are located in the project's own tag list staged in the snapshot (sources/stacks-chow/tags.txt), at the verbatim lines carrying their canonical tags.\par\medskip
\noindent\textit{Allowed standard prerequisites (editorial).}\begin{itemize}
\item Finiteness of presentation over a scheme theoretically dense open, openness of a flat and locally finitely presented morphism, and dimension zero fibres. Same chapter flat.tex, labels lemma-flat-finite-type-finitely-presented-over-dense-open-X (proof line 3084) and lemma-flat-finite-presentation-dimension-over-dense-open (proof line 3094); Morphisms, label morphisms-lemma-fppf-open (proof lines 3087 and 3104). The first upgrades local finite type to local finite presentation, the second shows all fibres have dimension zero whence the morphism is quasi-finite, and openness of the image lets the base be replaced by the image of the morphism.
\item The etale splitting of the quasi-finite part over an elementary etale neighbourhood. More on Morphisms, label more-morphisms-lemma-etale-splits-off-quasi-finite-part, invoked at proof line 3101; Morphisms, label morphisms-lemma-locally-quasi-finite-rel-dimension-0 (proof line 3096). Produces the elementary etale neighbourhood over which the base change splits into a finite part V and a part W empty over the closed fibre, which is the key representation of the morphism.
\item Scheme theoretic density under flat base change, descent of finiteness, and finite flat morphisms being finite locally free. Morphisms, labels morphisms-lemma-flat-morphism-scheme-theoretically-dense-open (proof line 3110), morphisms-lemma-finite-flat (proof line 3118) and morphisms-lemma-finite-locally-free (proof line 3121); Descent, label descent-lemma-descending-property-finite (proof line 3116). Shows the inverse image of the dense open remains scheme theoretically dense so that the empty part W must vanish, lets finiteness descend back to the base, and upgrades finite flat to finite locally free so that the degree can be compared over the dense open.
\end{itemize}\medskip
\ifdefined\hypertarget\hypertarget{sourcepointer-1}{}\fi
\noindent\textit{Source reference.} Original lemma 02LA: sources/stacks-chow/tags.txt, source lines 3401--3401. Navigation pointer for \textbackslash{}ref\{descent-lemma-descending-property-finite\} used inside the counted span (line 3116); the argument that consumes it is retained in the span.\par\medskip
\ifdefined\hypertarget\hypertarget{sourcepointer-2}{}\fi
\noindent\textit{Source reference.} Original lemma 081L: sources/stacks-flat/pinned-original-source.tex, source lines 3020--3020. Navigation pointer for \textbackslash{}ref\{lemma-flat-finite-presentation-dimension-over-dense-open\} used inside the counted span (line 3094); the argument that consumes it is retained in the span.\par\medskip
\ifdefined\hypertarget\hypertarget{sourcepointer-3}{}\fi
\noindent\textit{Source reference.} Original lemma 081P: sources/stacks-flat/pinned-original-source.tex, source lines 3000--3000. Navigation pointer for \textbackslash{}ref\{lemma-flat-finite-type-finitely-presented-over-dense-open-X\} used inside the counted span (line 3084); the argument that consumes it is retained in the span.\par\medskip
\ifdefined\hypertarget\hypertarget{sourcepointer-4}{}\fi
\noindent\textit{Source reference.} Original lemma 02LP: sources/stacks-chow/tags.txt, source lines 3416--3416. Navigation pointer for \textbackslash{}ref\{more-morphisms-lemma-etale-splits-off-quasi-finite-part\} used inside the counted span (line 3101); the argument that consumes it is retained in the span.\par\medskip
\ifdefined\hypertarget\hypertarget{sourcepointer-5}{}\fi
\noindent\textit{Source reference.} Original lemma 02KB: sources/stacks-chow/tags.txt, source lines 3366--3366. Navigation pointer for \textbackslash{}ref\{morphisms-lemma-finite-flat\} used inside the counted span (line 3118); the argument that consumes it is retained in the span.\par\medskip
\ifdefined\hypertarget\hypertarget{sourcepointer-6}{}\fi
\noindent\textit{Source reference.} Original lemma 04MH: sources/stacks-chow/tags.txt, source lines 6043--6043. Navigation pointer for \textbackslash{}ref\{morphisms-lemma-finite-locally-free\} used inside the counted span (line 3121); the argument that consumes it is retained in the span.\par\medskip
\ifdefined\hypertarget\hypertarget{sourcepointer-7}{}\fi
\noindent\textit{Source reference.} Original lemma 081H: sources/stacks-chow/tags.txt, source lines 10217--10217. Navigation pointer for \textbackslash{}ref\{morphisms-lemma-flat-morphism-scheme-theoretically-dense-open\} used inside the counted span (line 3110); the argument that consumes it is retained in the span.\par\medskip
\ifdefined\hypertarget\hypertarget{sourcepointer-8}{}\fi
\noindent\textit{Source reference.} Original lemma 01UA: sources/stacks-chow/tags.txt, source lines 2403--2403. Navigation pointer for \textbackslash{}ref\{morphisms-lemma-fppf-open\} used inside the counted span (line 3087); the argument that consumes it is retained in the span.\par\medskip
\ifdefined\hypertarget\hypertarget{sourcepointer-9}{}\fi
\noindent\textit{Source reference.} Original lemma 0397: sources/stacks-chow/tags.txt, source lines 4262--4262. Navigation pointer for \textbackslash{}ref\{morphisms-lemma-locally-quasi-finite-rel-dimension-0\} used inside the counted span (line 3096); the argument that consumes it is retained in the span.\par\medskip
\begin{lemma}
\label{lemma-zariski}
Let $f : X \to S$ be a morphism of schemes and $U \subset S$ an open. If
\begin{enumerate}
\item $f$ is separated, locally of finite type, and flat,
\item $f^{-1}(U) \to U$ is an isomorphism, and
\item $U \subset S$ is retrocompact and scheme theoretically dense,
\end{enumerate}
then $f$ is an open immersion.
\end{lemma}

\begin{proof}
By Lemma \ref{lemma-flat-finite-type-finitely-presented-over-dense-open-X}
the morphism $f$ is locally of finite presentation.
The image $f(X) \subset S$ is open
(Morphisms, Lemma \ref{morphisms-lemma-fppf-open})
hence we may replace $S$ by $f(X)$. Thus we have to prove that
$f$ is an isomorphism. We may assume $S$ is affine. We can reduce
to the case that $X$ is quasi-compact because it suffices to show
that any quasi-compact open $X' \subset X$ whose image is $S$
maps isomorphically to $S$. Thus we may assume $f$ is quasi-compact.
All the fibers of $f$ have dimension $0$, see
Lemma \ref{lemma-flat-finite-presentation-dimension-over-dense-open}.
Hence $f$ is quasi-finite, see
Morphisms, Lemma \ref{morphisms-lemma-locally-quasi-finite-rel-dimension-0}.
Let $s \in S$. Choose an elementary \'etale
neighbourhood $g : (T, t) \to (S, s)$ such that $X \times_S T = V \amalg W$
with $V \to T$ finite and $W_t = \emptyset$, see
More on Morphisms, Lemma
\ref{more-morphisms-lemma-etale-splits-off-quasi-finite-part}.
Denote $\pi : V \amalg W \to T$ the given morphism. Since $\pi$ is
flat and locally of finite presentation, we see that $\pi(V)$ is
open in $T$ (Morphisms, Lemma \ref{morphisms-lemma-fppf-open}).
After shrinking $T$ we may assume that $T = \pi(V)$.
Since $f$ is an isomorphism over $U$ we see that $\pi$ is an
isomorphism over $g^{-1}U$. Since $\pi(V) = T$ this implies
that $\pi^{-1}g^{-1}U$ is contained in $V$.
By Morphisms, Lemma
\ref{morphisms-lemma-flat-morphism-scheme-theoretically-dense-open}
we see that $\pi^{-1}g^{-1}U \subset V \amalg W$ is scheme theoretically
dense. Hence we deduce that $W = \emptyset$.
Thus $X \times_S T = V$ is finite over $T$.
This implies that $f$ is finite (after replacing $S$ by an
open neighbourhood of $s$), for example by
Descent, Lemma \ref{descent-lemma-descending-property-finite}.
Then $f$ is finite locally free
(Morphisms, Lemma \ref{morphisms-lemma-finite-flat})
and after shrinking $S$ to a smaller open neighbourhood of $s$
we see that $f$ is finite locally free of some degree $d$
(Morphisms, Lemma \ref{morphisms-lemma-finite-locally-free}).
But $d = 1$ as is clear from the fact that the degree is $1$ over
the dense open $U$. Hence $f$ is an isomorphism.
\end{proof}

\end{document}
